\begin{comment}
- in 1.1 we don't need too much detail about the el baf model yet, instead of saying it has KG, we can just say iot employs open-weight LLMs
- in 1.1 we also can't mention ontology without having defined it, we need to add a small sentence/paragraph explaining briefly what an ontology is and the reasoning behind its usage for our thesis/approach (in section 1.3)
in this we need to explain the reasoning, reference the most relevant paper.
- in the last sentence we can't say "ontology-based multi agent pipeline... ect." ; we can just explain the reasoning: (1) we are extending the el baff model with better context engineering, and our approach validates how we try to improve this context
- adding a small graph showing where exactly our approahc sits (sort of a miniature graph of the one in chapter 4) shoowing where we are working on (mostly in between retrieval and generation) and what we are not doing (working on retrieval or improving the underlying generation model).

- remove mistral small 24B model name in 1.3
- we need tio think about wether all these RQs are warranted and needed? are some not scoped or formulated correctly?
- we need to add a mini results section between 1.3 and 1.4: what's the takeaway from this thesis, how did the models perform in comparison? (part of the full thesis story). we can give the scores for the 4 criteria directly matching the 4 failure modes from 1.2

- in this same mini results section we can add a small part saying how (when given full text (when available), quantitive precision goes up at the cost of other metric performance, probably due to the longer contexts and difficulty to keep the numbers context aligned), here are the scores: set	set_label	scientific_accuracy_mean	relevance_mean	completeness_mean	directness_mean	quantitative_precision_mean	response_structure_mean	scope_awareness_mean	helpfulness_mean	overall_mean	scientific_accuracy_sd	relevance_sd	completeness_sd	directness_sd	quantitative_precision_sd	response_structure_sd	scope_awareness_sd	helpfulness_sd	overall_sd
set1	DLR 2-Step RAG (Large)	6.6571428571428575	6.614285714285714	6.571428571428571	4.585714285714285	4.071428571428571	5.9	4.542857142857143	5.128571428571429	5.508928571428571	0.5617187466255323	0.6207255488856065	0.627195935255346	0.6701256023920937	1.1956615847549839	0.6173810980989417	0.7553809728068843	0.7405722145536815	0.5406155193402727
set2	Ontology-Enhanced RAG (Small, abstracts)	7.828571428571428	7.957142857142857	7.457142857142857	7.928571428571429	5.928571428571429	7.671428571428572	7.628571428571429	7.414285714285715	7.476785714285715	0.6587521926686511	0.5757342957441968	0.7553809728068843	0.6213922789906823	1.5258300527570168	0.7934775518304215	0.8370311332396319	0.8596065352503636	0.6314137781306728
set3	No-Ontology RAG	7.185714285714286	7.557142857142857	6.857142857142857	7.114285714285714	5.328571428571428	7.328571428571428	7.071428571428571	6.557142857142857	6.875	0.8391309711537779	0.9268258796452138	1.0113230577427537	0.6924616275217437	1.5296246444128152	0.756065877266026	0.6442933118041897	1.0852865017543842	0.7402751153299776
set4	Zero-Shot	7	7.571428571428571	6.914285714285715	6.671428571428572	4.014285714285714	7.3	5.3428571428571425	5.414285714285715	6.2785714285714285	0.6370220572706061	0.6930593502290829	0.6966349166093654	0.6532289602690327	1.2216982376392838	0.5477225575051661	0.866204686552902	0.8425781199382986	0.5603002189931097
set5	Ontology+Full-Text (Small)	7.228571428571429	7.557142857142857	7.214285714285714	7.085714285714285	6.828571428571428	7.714285714285714	7.042857142857143	6.914285714285715	7.198214285714286	0.6409103256686182	0.8277035506157161	1.0340977603949397	0.6966349166093654	1.6415919153562943	0.8190283593479222	0.9236931902589999	1.2481041523668126	0.7788802049776359
set6	Ontology+Full-Text (Large)	7.114285714285714	7.642857142857143	7.585714285714285	7.3428571428571425	7.814285714285714	8.014285714285714	7.185714285714286	7.457142857142857	7.519642857142857	0.8434377057975749	0.8171302609462198	1.0284765311640707	0.6344166391652172	0.8562279451150022	0.6701256023920937	0.7281673868494946	0.7553809728068843	0.6264966552234825


\end{comment}


\chapter{Introduction}
\label{chap:introduction}

\section{Context and Motivation}
\label{sec:context}

Large language models (LLMs) hallucinate, and in scientific domains these errors carry severe consequences. The most direct mitigation is to ground generation in retrieved evidence, the core idea behind Retrieval-Augmented Generation (RAG). A RAG system adds documents such as papers, articles, and web pages to the model's input, retrieved by term-based methods, semantic embeddings, or both. This approach has become central to reliable AI in scientific domains, where answers need correct terminology, quantitative precision, and explicit source attribution~\citep{lewis2020retrieval}.

Earth observation (EO) is particularly sensitive to all three. It spans remote sensing, climate monitoring, and environmental assessment, and it relies on specialised vocabulary, quantitative accuracy, and strict temporal and spatial context~\citep{aschbacher2020earth}. Each measurement or study only makes sense in context: which sensors were used, what was measured, when and where, and under which conditions.

\citet{elbaff2025earthobservationrag} propose a knowledge-graph RAG pipeline for open-source EO question answering, combining lexical and semantic retrieval with a structured generation step. We assess why larger proprietary language models such as GPT-5.1~\citep{openai2024gpt4} and Claude Sonnet~4.5~\citep{anthropic2024claude} handle EO queries well while this open-source model underperforms. We use this model and the data generated by its architecture as the \emph{baseline} throughout this thesis. All proposed extensions are compared against it.

Standard RAG approaches, including the baseline, retrieve documents using lexical and semantic search, without modelling the semantic structure of the query. Even when relevant documents are retrieved, the assembled context fails to support generation, critical details are lost, and irrelevant passages dilute the answer to the point where it cannot support reliable scientific reasoning. This thesis addresses these failures by designing, implementing, and evaluating an
ontology-based multi-agent pipeline that imposes semantic structure on context construction and generation.
\section{Problem Statement}
\label{sec:problem}

After systematic manual analysis of baseline answers, the baseline~\citep{elbaff2025earthobservationrag} proved to fail on specialised and complex queries due to a lack of deeper contextual awareness. The shortcomings identified below are expanded with detailed examples in Chapter~\ref{chap:preliminaries}.

\paragraph{Key Shortcomings}\label{para:shortcomings}

We evaluated how well the baseline handled questions of varying complexity and type, focusing on domain-specific queries that combine multiple concepts and require both reasoning and specialised data. We compared the baseline against Claude Sonnet~4.5~\citep{anthropic2024claude} and GPT-5.1~\citep{openai2024gpt4} as proprietary reference systems. After iterative prompt refinement, we observe four main shortcomings:

\subparagraph{1. Delayed direct answers}

A direct question warrants a direct answer. The baseline opens with several sentences of background and context only superficially related to the query, burying the actual answer. Proprietary models consistently lead with the answer itself, offering context only to support it.

\subparagraph{2. Missing quantitative precision}

Scientific answers should extract and report the specific numbers that matter: thresholds, metrics, and the conditions under which they were obtained. The baseline mentions relevant papers without extracting these quantitative details, even when the cited sources contain them explicitly.

\subparagraph{3. Lack of useful structure}

Not every question is best answered in prose; tables and lists are often more appropriate for comparisons and multi-part queries. The baseline defaults to long paragraphs where key points are scattered and mixed with background, regardless of the question type. Proprietary models are able to adapt their response format to the nature of the question.

\subparagraph{4. Context misalignment}

Domain-specific questions combine multiple contexts and dimensions, and the answer must make clear which sources apply to which setting. The baseline mixes contexts, creating a risk of serious misunderstanding when experts assume the reported information is specific to their query. Proprietary models often explicitly flag when a cited source addresses a different setting.

\section{Approach and Research Questions}
\label{sec:research-questions}

The four shortcomings identified above share a root cause: the baseline retrieves and generates without any explicit representation of what the query is actually asking across its semantic dimensions. To address this, we propose a three-agent pipeline in which an \emph{ontology agent} first parses each query into a structured, query-specific ontology, a lightweight set of attribute-value pairs capturing dimensions such as temporal scope, spatial scope, sensor or method, and biophysical variable. A \emph{refinement agent} uses this ontology to filter retrieved documents and synthesise a curated \emph{context}. A \emph{generation agent} then drafts and iteratively refines the final answer against this curated context. All agents run on Mistral Small 24B, the same open-source model used in the baseline. The following research questions structure the investigation of this approach.

\textbf{Main Research Question:} Can a multi-agent architecture with dynamic ontology generation improve answer quality in domain-specific RAG systems for Earth observation queries?

\textbf{Sub-Research Question 1:} How can query-specific ontologies capture the semantic structure of Earth observation questions to guide document retrieval and context refinement?

\textbf{Sub-Research Question 2:} To what extent can small language models handle complex Earth observation queries within a multi-agent RAG pipeline?

\textbf{Sub-Research Question 3:} How reliable are LLM-as-a-judge evaluation frameworks for assessing RAG systems in specialised scientific domains, and what are their key limitations?

\textbf{Sub-Research Question 4:} Does ontology-based semantic grounding provide significant performance benefits for small-model RAG systems, and to what extent can such an approach generalise to other knowledge-intensive domains beyond Earth observation?

\section{Thesis Contributions}
\label{sec:contribution}

This thesis makes the following contributions:

\begin{enumerate}
\item \textbf{A structured failure taxonomy of open-source RAG for Earth observation.}
Through qualitative analysis of the baseline, we identify four recurring failure
categories: delayed directness, missing quantitative precision, poor response structure,
and context misalignment. These categories reveal a systematic gap between open-source
and proprietary model behaviour on domain-specific queries.

\item \textbf{A novel ontology-based multi-agent pipeline for domain-specific RAG.}
We present a three-agent architecture, comprising an ontology agent, a refinement agent,
and a generation agent, in which each stage enforces explicit semantic reasoning about
the query. The pipeline generates a compact, query-specific ontology covering five
dimensions (temporal scope, spatial scope, application domain, sensor/method, and
biophysical variable) for every query, rather than relying on a static domain taxonomy,
directly targeting the root causes of the identified failure modes.

\item \textbf{A tailored multi-criteria evaluation framework for scientific RAG.}
We develop an eight-criterion LLM-as-a-judge framework using Claude
Sonnet~4.6~\citep{anthropic2024claude} as judge on a 1--10 scale, specifically designed
to quantify the attributes reflecting the four shortcomings. The framework addresses
score saturation and context-length biases identified in an initial evaluation
configuration (see Chapter~\ref{chap:evaluation}).

\item \textbf{Empirical evidence that ontology-based semantic grounding significantly
improves answer quality.}
Evaluated across 70 EO questions~\citep{elbaff2025earthobservationrag}, our pipeline
outperforms the baseline by 1.97 points overall and ranks first in all eight criteria,
with all pairwise comparisons significant at $p < 0.001$ and effect sizes
$r \geq 0.77$.

\item \textbf{A qualitative error analysis revealing the limits of abstract-only
retrieval.}
Manual inspection of failure cases shows that quantitative precision, the weakest
criterion across all variants, is systematically limited by the use of document
abstracts, which rarely contain exact measurements, pointing to full-text retrieval as
the most impactful direction for future work.
\end{enumerate}

The principles of query-adaptive ontology generation and semantic document refinement
are not exclusive to Earth observation: any domain where queries combine multiple
semantic dimensions stands to benefit from this approach.

\section{Thesis Structure}
\label{sec:structure}

Chapter~\ref{chap:background} covers background on RAG systems and Earth observation, and reviews related work on ontology-grounded retrieval, multi-agent systems, and LLM-based evaluation. Chapter~\ref{chap:preliminaries} deepens the analysis of the baseline model's shortcomings with detailed examples. Chapter~\ref{chap:approach} presents the three-agent architecture, specifying each agent's role and showing how the design addresses each identified shortcoming. Chapter~\ref{chap:evaluation} describes the evaluation framework: the dataset of 70~EO questions, four pipeline variants, and the multi-criteria LLM-as-a-judge methodology. Chapter~\ref{chap:results} reports quantitative results, compares the four variants, and identifies failure modes through detailed examples. Chapter~\ref{chap:discussion} interprets the results and discusses limitations. Finally, Chapter~\ref{chap:conclusion} reflects on possible future work and concludes by revisiting the research questions.
